\section{Conclusion}
\label{sec:conclusion}

Block-diffusion language models have a staircase cost shape that is
structurally different from the linear cost of autoregressive decoding.
The shapes cross at $L^{*} = \crossoverL$ output tokens, above which
diffusion is faster --- but only when the comparison is fair.

Against a MoE-matched AR model (Qwen3-30B-A3B, 3\,B active of 30\,B),
the diffusion model's advantage is modest: $1.4\times$ at $L{=}1024$,
$2.4\times$ at $L{=}2048$. Below 368 tokens, AR is faster. The six-fold
improvement in per-token cost from dense-to-MoE AR --- 6.3 to
1.0\,ms/token --- is the single largest factor shrinking the niche.

The advantage does not survive deployment. The AR model batches
$77\times$ better, has a comparable footprint (60\,GB vs 52\,GB), and
a $3.21\times$ tighter tail. Under tenant-isolated batch-1, sixteen
copies of the quality-similar 3\,GB dense AR model fit in the same GPU,
making it $9\times$ cheaper per event.

What survives is an extremely narrow regime: output above $\sim$1000
tokens, batch-1 by requirement, schema validation, and a data-centre
footprint. Outside that regime, the line beats the staircase.
